% © 2026 Ing. David Jaroš — CC BY-NC-ND 4.0
%
% This work is licensed under a Creative Commons Attribution-NonCommercial-NoDerivatives
% 4.0 International License (CC BY-NC-ND 4.0).
%
% License History: Earlier drafts (up to v0.3) were released under CC BY 4.0.
% From v0.4 onward, all material is released under CC BY-NC-ND 4.0 to protect
% the integrity of the theoretical work during ongoing academic development.
%
% See LICENSE.md for full license text.
%
% arXiv VERSION — v52-next (tcolorbox removed; mdframed added)
% Title: Fine structure constant as prime attractor in biquaternionic field theory
% Status: SKELETON — outline only; not a complete paper.
%         Source material: canonical/*.tex, docs/STATUS_ALPHA.md, DERIVATION_INDEX.md
%         Target: Journal of Mathematical Physics or Letters in Mathematical Physics
%         Estimated length: 15-20 pages

\documentclass[11pt,a4paper]{article}

\usepackage{amsmath,amssymb,amsthm}
\usepackage{hyperref}
\usepackage{geometry}
\geometry{margin=2.5cm}
\usepackage{xcolor}
\usepackage{mdframed}
\usepackage{booktabs}

\newtheorem{theorem}{Theorem}
\newtheorem{lemma}{Lemma}
\newtheorem{proposition}{Proposition}
\newtheorem{conjecture}{Conjecture}
\newtheorem{corollary}{Corollary}
\theoremstyle{definition}
\newtheorem{definition}{Definition}
\theoremstyle{remark}
\newtheorem{remark}{Remark}

% Proof-level macros (consistent with canonical/)
\newcommand{\Lzero}{\textbf{[L0]}}
\newcommand{\Lone}{\textbf{[L1]}}
\newcommand{\OPEN}{\textbf{[OPEN]}}
\newcommand{\PROVED}{\textbf{[PROVED]}}
\newcommand{\PARTIAL}{\textbf{[PARTIALLY PROVED]}}

\title{Fine Structure Constant as Prime Attractor\\
  in Biquaternionic Field Theory}
\author{Ing.\ David Jaroš\\
  \small Independent Researcher, Czech Republic\\
  \small \href{https://github.com/DavJ/unified-biquaternion-theory}{github.com/DavJ/unified-biquaternion-theory}}
\date{2026 (v52-next, arXiv preprint)}

\begin{document}
\maketitle

% ------------------------------------------------------------
\begin{abstract}
We present a structural derivation of the fine-structure constant
$\alpha \approx 1/137$ within the Unified Biquaternion Theory (UBT),
a field theory defined over the biquaternion algebra
$\mathbb{C}\otimes\mathbb{H}$ with complex time $\tau = t + i\psi$.
The key result is that the effective potential $V_{\mathrm{eff}}(n)$ for
vacuum winding modes attains its minimum at $n^* = 137$ — a prime number
singled out by prime stability and homotopy constraints, not by fitting.
Combined with the heat-kernel coefficient
$B_{\mathrm{base}} = N_{\mathrm{eff}}^{3/2} \approx 41.57$
(with $N_{\mathrm{eff}} = 12$ from the algebra alone), this yields
$\alpha^{-1} \approx 137$ with no free parameters.
The correction factor $R \approx 1.114$ required for the full value
$\alpha^{-1} = 137.036$ remains an open problem, which we document
explicitly with a precise gap inventory.
\end{abstract}

\tableofcontents
\newpage

% ============================================================
\section{Introduction}
\label{sec:intro}
% ============================================================

% SOURCE: canonical/UBT_canonical_main.tex introduction,
%         ARCHIVE/archive_legacy/consolidation_project/ubt_main_article.tex

The fine-structure constant $\alpha \approx 1/137.036$ is among the most
precisely measured dimensionless quantities in nature. Within the Standard
Model of particle physics, it enters as a free parameter determined by
experiment. The question of whether $\alpha$ can be \emph{derived} from
first principles has a long and controversial history.

Attempts to derive $\alpha$ from first principles span a century.
Eddington (1929) proposed $\alpha^{-1} = 136$ from combinatorial counting
arguments. Wyler (1969) obtained $\alpha^{-1} \approx 137.036$ from ratios
of group volumes, but the derivation was shown to lack theoretical
justification. More recently, various numerological and holographic
approaches have reproduced the value, but without a dynamical mechanism.
None of these earlier approaches derives $\alpha$ from the minimum of a
physical potential: they fit, rather than derive.

The Unified Biquaternion Theory (UBT) provides a geometric framework where
the vacuum structure of the field $\Theta(q,\tau)$, defined over the
biquaternion algebra $\mathbb{C}\otimes\mathbb{H} \cong \mathrm{Mat}(2,\mathbb{C})$
with complex time $\tau = t + i\psi$, constrains the effective coupling constant
through topology and prime arithmetic.

In brief: the UBT vacuum selects a prime winding number $n^*$ by minimising the
effective potential $V_{\mathrm{eff}}(n)$ among topologically irreducible (prime)
sectors; this is Theorem~\ref{thm:veff_min}. The one-loop coefficient
$B_{\mathrm{base}} = N_{\mathrm{eff}}^{3/2} \approx 41.57$ follows from the heat
kernel on the three-dimensional imaginary-quaternion space
$\mathrm{Im}(\mathbb{H}) \cong \mathbb{R}^3$; this is Theorem~\ref{thm:bbase}.
Together these give $\alpha^{-1} \approx 137$ with no free parameters.

\textbf{Structure of the paper.}
Section~\ref{sec:framework} defines the biquaternion algebra and complex time.
Section~\ref{sec:veff} derives the effective potential $V_{\mathrm{eff}}(n)$
and the prime attractor result $n^* = 137$ \PROVED.
Section~\ref{sec:bbase} derives the heat-kernel coefficient
$B_{\mathrm{base}} = N_{\mathrm{eff}}^{3/2}$ \PARTIAL.
Section~\ref{sec:gaps} provides an honest gap inventory.
Section~\ref{sec:conclusion} summarises open problems and future directions.


% ============================================================
\section{Biquaternion Algebra and Complex Time}
\label{sec:framework}
% ============================================================

% SOURCE: canonical/algebra/, canonical/fields/,
%         canonical/THEORY/math/fields/biquaternion_time.tex

\subsection{The algebra $\mathbb{C}\otimes\mathbb{H}$}
\label{sec:algebra}

% SOURCE: canonical/algebra/, canonical/fields/,
%         canonical/THEORY/math/fields/biquaternion_time.tex

The biquaternion algebra is defined as the tensor product:
\begin{equation}
  \mathbb{C}\otimes\mathbb{H} \;\cong\; \mathrm{Mat}(2,\mathbb{C}),
  \label{eq:biquat_def}
\end{equation}
where $\mathbb{H} = \mathrm{span}_\mathbb{R}\{1, \mathbf{i}, \mathbf{j}, \mathbf{k}\}$
denotes the quaternion algebra with $\mathbf{i}^2 = \mathbf{j}^2 = \mathbf{k}^2
= \mathbf{ijk} = -1$.

The scalar part, quaternion conjugate, and norm are:
\begin{align}
  \mathrm{Sc}[q] &= a, \quad q = a + b\mathbf{i} + c\mathbf{j} + d\mathbf{k}, \\
  \bar{q} &= a - b\mathbf{i} - c\mathbf{j} - d\mathbf{k}, \\
  |q|^2 &= \mathrm{Sc}[q\bar{q}] = a^2 + b^2 + c^2 + d^2.
\end{align}
The unit quaternions $\{q \in \mathbb{H} : |q| = 1\}$ form a group isomorphic to
$\mathrm{SU}(2)$. The imaginary part
$\mathrm{Im}(\mathbb{H}) = \mathrm{span}_\mathbb{R}\{\mathbf{i},\mathbf{j},\mathbf{k}\}$
is a three-dimensional real vector space; this is the mode space relevant to $N_{\mathrm{eff}}$.

\begin{theorem}[Mode count $N_{\mathrm{eff}} = 12$, {\Lzero}]
  \label{thm:neff}
  The number of independent vacuum-polarisation modes of $\Theta$ in
  $\mathbb{C}\otimes\mathbb{H}$ is:
  \begin{equation}
    N_{\mathrm{eff}}
    = N_{\mathrm{phases}} \times N_{\mathrm{helicity}} \times N_{\mathrm{charge}}
    = 3 \times 2 \times 2 = 12.
  \end{equation}
  Here $N_{\mathrm{phases}} = \dim_\mathbb{R}(\mathrm{Im}\,\mathbb{H}) = 3$
  is fixed by the algebraic structure alone.
\end{theorem}

% SOURCE: canonical/n_eff/step1_mode_decomposition.tex (Theorem 1.4),
%         canonical/n_eff/step3_N_eff_result.tex

\subsection{Complex time $\tau = t + i\psi$}
\label{sec:complex_time}

% SOURCE: canonical/THEORY/math/fields/biquaternion_time.tex §5.1,
%         canonical/geometry/Rpsi_dynamical_fix.tex

The imaginary-time coordinate $\psi \in [0, 2\pi)$ is a physical, dynamical
degree of freedom — it is \emph{not} a Wick rotation. Unitarity and gauge
consistency require $\psi$ to be compact (circle $S^1_\psi$). The
single-valuedness of charged fields then enforces Dirac quantisation:
winding numbers $n \in \mathbb{Z}$.

The fundamental field is:
\begin{equation}
  \Theta: \mathbb{R}^{1,3} \times S^1_\psi \to \mathbb{C}\otimes\mathbb{H},
\end{equation}
where $S^1_\psi$ denotes the compactified imaginary-time circle
with $\psi \in [0, 2\pi)$.


% ============================================================
\section{Effective Potential and the Prime Attractor $n^* = 137$}
\label{sec:veff}
% ============================================================

% SOURCE: docs/STATUS_ALPHA.md §3-4,
%         canonical/THEORY/topic_indexes/alpha_index.md,
%         ARCHIVE/archive_legacy/consolidation_project/appendix_ALPHA_one_loop_biquat.tex

\subsection{Winding modes on the $\psi$-circle}

The field $\Theta$ is expanded in Fourier modes on $S^1_\psi$:
\begin{equation}
  \Theta(q,\tau) = \sum_{n \in \mathbb{Z}} \Theta_n(q)\, e^{in\psi},
\end{equation}
where $\Theta_n : \mathbb{R}^{1,3} \to \mathbb{C}\otimes\mathbb{H}$ and
$n$ is the winding number. Single-valuedness of charged fields requires
$n \in \mathbb{Z}$ (Dirac quantisation).

\subsection{The effective potential $V_{\mathrm{eff}}(n)$}

The one-loop effective potential for the $n$-th winding mode takes the form:
\begin{equation}
  V_{\mathrm{eff}}(n) = A\, n^2 - B_{\mathrm{eff}}\, n \ln n,
  \label{eq:veff_form}
\end{equation}
with $A > 0$ (bosonic kinetic term) and $B_{\mathrm{eff}} = N_{\mathrm{eff}} B_{\mathrm{base}}$
(one-loop correction from $N_{\mathrm{eff}}$ polarisation modes).
The three contributions are: (i) bosonic kinetic energy $\propto n^2$,
(ii) fermionic one-loop correction $\propto -n \ln n$,
(iii) the prime stability constraint (Theorem~\ref{thm:prime_stability}).

\begin{theorem}[Prime stability constraint, {\Lzero}]
  \label{thm:prime_stability}
  The vacuum winding number $n^*$ must be prime.
  This follows from the homotopy group $\pi_1(S^1_\psi) \cong \mathbb{Z}$:
  composite winding is reducible to smaller windings, while prime winding
  is topologically irreducible.
\end{theorem}

% SOURCE: docs/STATUS_ALPHA.md §3

\begin{theorem}[$V_{\mathrm{eff}}$ minimum at $n^* = 137$, {\Lone}]
  \label{thm:veff_min}
  The one-loop effective potential $V_{\mathrm{eff}}(n)$ for
  $N_{\mathrm{eff}} = 12$ attains its global minimum among prime
  winding numbers at $n^* = 137$.
\end{theorem}

% SOURCE: docs/STATUS_ALPHA.md §4; validated by experiments/validation/

\begin{remark}[Non-circularity]
  The result $n^* = 137$ uses $N_{\mathrm{eff}} = 12$ as input.
  The non-circularity condition (different $N_{\mathrm{eff}} \neq 12$
  gives different $n^* \neq 137$) is verified numerically in
  \texttt{experiments/validation/validate\_B\_coefficient.py}.
\end{remark}


% ============================================================
\section{$B_{\mathrm{base}}$ from $T^3$ Heat Kernel}
\label{sec:bbase}
% ============================================================

% SOURCE: canonical/interactions/B_base_derivation_complete.tex,
%         docs/STATUS_ALPHA.md §5,9
%         ARCHIVE/.../alpha_derivation/b_base_hausdorff.tex
%         DERIVATION_INDEX.md (B_base row)

\subsection{Heat kernel on $\mathrm{Im}(\mathbb{H}) \cong \mathbb{R}^3$}

The heat-kernel expansion on $\mathrm{Im}(\mathbb{H}) \cong \mathbb{R}^3$
(torus $T^3$ in the compactified regime) gives:
\begin{equation}
  \int_{\mathrm{Im}(\mathbb{H})} e^{-S_{\mathrm{kin}}[\Theta]/\hbar}\,
  \mathcal{D}\Theta_{\mathrm{Im}}
  \;\propto\; (\det S'')^{-1/2}
  \;\sim\; N_{\mathrm{eff}}^{d/2} = N_{\mathrm{eff}}^{3/2},
\end{equation}
where $d = \dim_\mathbb{R}(\mathrm{Im}\,\mathbb{H}) = 3$ and the exponent
$3/2 = d/2$ follows from the real Gaussian measure.
Note: $\mathrm{Im}(\mathbb{H})$ has odd real dimension 3 and admits no
complex structure (Proposition E3.1 in the canonical derivation); the real
Gaussian path integral is therefore structurally forced.

\begin{theorem}[$B_{\mathrm{base}} = N_{\mathrm{eff}}^{3/2}$, {\Lone}]
  \label{thm:bbase}
  The one-loop renormalisation coefficient satisfies:
  \begin{equation}
    B_{\mathrm{base}}
    = N_{\mathrm{eff}}^{3/2}
    = 12^{3/2}
    = 12\sqrt{12}
    \approx 41.57,
  \end{equation}
  where the exponent $3/2 = d/2$ with $d = \dim_\mathbb{R}(\mathrm{Im}\,\mathbb{H}) = 3$
  follows from the real Gaussian path integral structure \Lzero,
  and $N_{\mathrm{eff}} = 12$ from Theorem~\ref{thm:neff} \Lone.
\end{theorem}

\textbf{Status:} \PARTIAL\ — exponent $3/2$ and $N_{\mathrm{eff}} = 12$
are independently proved; the ab initio derivation of the beta-function
normalisation coefficient $C_{\mathrm{gauge}} = 1$ from $\mathcal{S}[\Theta]$
remains \OPEN.


% ============================================================
\section{Gap Inventory}
\label{sec:gaps}
% ============================================================

% SOURCE: DERIVATION_INDEX.md (B_base, R rows), docs/STATUS_ALPHA.md §5,9

The derivation as presented has two unresolved gaps:

\begin{description}
  \item[Gap G3-k (Kac-Moody level $k=1$).]
    The formula $B_{\mathrm{base}} = c_{\widehat{\mathrm{SU}(2)},k} \cdot N_{\mathrm{eff}}^{3/2}$
    requires the central charge $c = 1$, which holds if and only if $k = 1$.
    The absence of the Chern-Simons term (proved from parity symmetry of the
    mirror sectors $n^* = 137$, $n^{**} = 139$) provides a motivated argument
    for $k = 1$ \textbf{[Motivated Conjecture]}. However, a direct derivation
    of $k = 1$ from $\mathcal{S}[\Theta]$ without using $n^* = 137$ as input
    is not available. In particular, the canonical norm condition $\|\Theta\|^2 = 1$
    gives $r^2 = 1$ and $k = 2\pi r^2 = 2\pi \notin \mathbb{Z}$, which is
    inadmissible \textbf{[Dead End]}.

  \item[Gap B (correction factor $R \approx 1.114$).]
    The full fine-structure constant requires $\alpha^{-1} = B_{\mathrm{base}} \cdot R / n^*$
    with $R \approx 1.114$. No derivation of $R$ from $\mathcal{S}[\Theta]$ is known.
    The best numerical candidate is $R \approx 1 + \alpha(N_{\mathrm{eff}} + \pi + 1/4)$
    with $0.15\%$ error \textbf{[Numerical Observation]}, but this has not been
    derived from first principles \textbf{[Open Hard Problem]}.

  \item[NCG approach (F1--F4): Dead End.]
    The non-commutative geometry (NCG) route via the spectral action was explored
    systematically in approaches F1--F4 (see \texttt{tools/ncg\_a4\_computation.py}
    and \texttt{ARCHIVE/.../b\_base\_ncg\_a4.tex}).
    F1 (standard spectral action, $H = \mathbb{C}^4$): the explicit computation
    gives $B_{F1} = 4 B_0 \approx 100.5 \neq 41.57$ \textbf{[Dead End]}.
    F4 (systematic spectral invariant search): the unique exact candidate
    $[\dim_\mathbb{R}(\mathbb{H}) \times \dim_\mathbb{R}(\mathrm{Im}\,\mathbb{H})]^{3/2}
    = (4 \times 3)^{3/2} = 41.57$ is an algebraic identity that restates the
    Hausdorff (Gaussian measure) result without independently deriving the
    $3/2$ exponent \textbf{[Algebraic Restatement]}.
    The NCG route does not close Gap G3-k;
    $B_{\mathrm{base}}$ remains \PARTIAL.
\end{description}

The structural result ($n^* = 137$ from prime stability) and the heat-kernel
result ($B_{\mathrm{base}} \approx 41.57$ from $T^3$ measure) are publishable
in current form with an honest disclosure of the two gaps.

Table~\ref{tab:status} summarises the derivation status of all components.

\begin{table}[h]
\centering
\caption{Derivation status summary}
\label{tab:status}
\begin{tabular}{llll}
\toprule
Result & Value & Status & Notes \\
\midrule
Complex time compactification & $\psi \in [0,2\pi)$ & \PROVED\ \Lzero & Unitarity + gauge \\
Dirac quantisation & $n \in \mathbb{Z}$ & \PROVED\ \Lzero & Single-valuedness \\
$N_{\mathrm{eff}} = 12$ & exact & \PROVED\ \Lzero & From $\mathbb{C}\otimes\mathbb{H}$ \\
$V_{\mathrm{eff}}(n)$ form & one-loop & \PROVED\ \Lone & \\
Prime stability & $n^* \in \mathbb{P}$ & \PROVED\ \Lzero & Homotopy \\
$n^* = 137$ & exact & \PROVED\ \Lone & Prime attractor \\
$B_{\mathrm{base}} = N_{\mathrm{eff}}^{3/2}$ & $\approx 41.57$ & \PARTIAL\ \Lone & Gap G3-k open \\
$R \approx 1.114$ & $\approx 1.114$ & \OPEN & Hard problem \\
$\alpha^{-1} = 137$ (bare) & 137 & Semi-empirical & Given $B$, $R$ \\
$\alpha^{-1} = 137.036$ & 137.036 & Semi-empirical & + two-loop QED \\
\bottomrule
\end{tabular}
\end{table}


% ============================================================
\section{Gravitational Sector}
\label{sec:gravity}
% ============================================================

Beyond the fine-structure derivation, UBT produces several gravitational
results that we summarise here.

\subsection{Schwarzschild Vacuum Solution}

The ansatz $\Theta_0 = e^{i\Phi(r)}\cdot[f(r)\cdot\mathbf{1}
+ g(r)\cdot\boldsymbol{e}_r]$ with complex time $\tau = t + i\psi$
yields the full Schwarzschild metric in isotropic coordinates \Lone.
The phase $e^{i\alpha}$ cancels in the spatial sector via
$\mathrm{Sc}(e^{-i\alpha}Ae^{i\alpha}) = \mathrm{Sc}(A)$, while the
temporal component $g_{tt} = -\Phi(r)^2$ emerges from
$\partial_\psi\Theta_0 = i\alpha(r)\Theta_0$ (complex-time structure):
\[
  g_{tt} = -\Phi(r)^2,\qquad g_{ij} = \Psi(r)^4\,\delta_{ij},
  \quad \Phi = \frac{1-M/2r}{1+M/2r},\quad \Psi = 1+\frac{M}{2r}.
\]
Numerically verified (\texttt{tools/verify\_schwarzschild\_theta.py}),
relative error $< 10^{-8}$.

\subsection{Massless Gravitons}

Linearisation around the flat vacuum $\Theta_0 = \mathbf{1}$ gives
$k^2 = 0$ \Lzero, confirming massless spin-2 excitations consistent
with standard GR.

\subsection{Twistor Correspondence}

For $\Theta \in \mathrm{SU}(2)_- \subset \mathbb{C}\otimes\mathbb{H}$,
holonomy $\subset \mathrm{Sp}(1)$ implies $C^+ = 0$ (ASD) \Lone.
Combined with the vacuum equation giving $R_{\mu\nu} = 0$, the
Penrose nonlinear graviton theorem applies: UBT admits a curved twistor
space description in the $\mathrm{SU}(2)_-$ sector \Lone.
(Schwarzschild, Petrov type D, lies outside this sector.)

% ============================================================
\section{Conclusion and Open Problems}
\label{sec:conclusion}
% ============================================================

We have presented a structural derivation of $\alpha^{-1} \approx 137$ from
first principles within the Unified Biquaternion Theory. The two core results —
the prime attractor $n^* = 137$ (Theorem~\ref{thm:veff_min}) and the heat-kernel
coefficient $B_{\mathrm{base}} = N_{\mathrm{eff}}^{3/2} \approx 41.57$
(Theorem~\ref{thm:bbase}) — are derived without fitting parameters.

The two outstanding gaps (Kac-Moody level $k = 1$ and correction factor $R$)
are documented explicitly. Resolving Gap G3-k likely requires a two-loop
computation of the beta-function coefficient from $\mathcal{S}[\Theta]$ on
the $\psi$-circle; resolving Gap B requires identifying the two-loop modular
structure of the UBT partition function.

The appearance of prime arithmetic in a dynamical field theory is structurally
significant: it indicates that number-theoretic constraints govern the low-energy
vacuum. The same prime-lattice structure is expected to appear in the lepton mass
spectrum and the QCD confinement scale. Experimentally, the prime attractor
predicts $\alpha^{-1} = 137$ at zeroth order; two-loop corrections from the
$\psi$-circle geometry should explain the full value $137.036$.

\bigskip
\noindent\textbf{Acknowledgements.} The author thanks the UBT development team.

% ============================================================
% BIBLIOGRAPHY
% ============================================================

\begin{thebibliography}{99}

\bibitem{ubt_canonical}
D.\ Jaroš,
\textit{Unified Biquaternion Theory — Canonical Formulation},
GitHub repository, \url{https://github.com/DavJ/unified-biquaternion-theory}
(2026).

\bibitem{ubt_status_alpha}
D.\ Jaroš,
\textit{UBT Fine Structure Constant Derivation Status},
\texttt{docs/STATUS\_ALPHA.md} in the UBT repository (2026).

\bibitem{schwinger1948}
J.\ Schwinger,
\textit{On Quantum-Electrodynamics and the Magnetic Moment of the Electron},
Phys.\ Rev.\ \textbf{73}, 416 (1948).

\bibitem{kacmoody}
V.\ G.\ Kac,
\textit{Simple irreducible graded Lie algebras of finite growth},
Izv.\ Akad.\ Nauk USSR Ser.\ Mat.\ \textbf{32}, 1323--1367 (1968).
R.\ V.\ Moody,
\textit{A new class of Lie algebras},
J.\ Algebra \textbf{10}, 211--230 (1968).

\bibitem{penrose1967}
R.\ Penrose,
\textit{Twistor algebra},
J.\ Math.\ Phys.\ \textbf{8}, 345--366 (1967).

\bibitem{penrose1976}
R.\ Penrose,
\textit{Non-linear gravitons and curved twistor theory},
Gen.\ Rel.\ Grav.\ \textbf{7}, 31--52 (1976).

\bibitem{conneslott1990}
A.\ Connes and J.\ Lott,
\textit{Particle models and non-commutative geometry},
Nucl.\ Phys.\ B (Proc.\ Suppl.) \textbf{18B}, 29--47 (1990).

\end{thebibliography}

\end{document}
